% ==================================================================
%  Обложка и страница содержания
%  Требует определённых в meta.tex: \subjectname, \subjectshort,
%  \lecturer, \subjectcolor, \coursemeta, \compiler, \compilergh
% ==================================================================

\usepackage{anyfontsize}
\usepackage{tocloft}

% ---------- содержание --------------------------------------------
\renewcommand{\cfttoctitlefont}{\sffamily\bfseries\Large\color{ink}}
\renewcommand{\cftsecfont}{\sffamily\color{ink}}
\renewcommand{\cftsecpagefont}{\sffamily\color{inkSoft}}
\renewcommand{\cftsecleader}{\color{ruleC}\cftdotfill{\cftdotsep}}
\renewcommand{\cftsecnumwidth}{2.2em}
\setlength{\cftbeforesecskip}{5pt}
\renewcommand{\cftaftertoctitleskip}{12pt}

% ---------- обложка -----------------------------------------------
\newcommand{\makecover}{%
  \begin{titlepage}
  \thispagestyle{empty}
  \begingroup\sffamily
  \begin{tikzpicture}[remember picture, overlay,
    % TikZ фиксирует \spaceskip по внешнему кеглю; без сброса межсловные
    % пробелы в крупном заголовке схлопываются
    every node/.append style={execute at begin node={\spaceskip=0pt\relax}}]
    \fill[\subjectcolor] (current page.south west) rectangle (current page.north east);

    % генеративный фон: диагональная сетка + дуги + несколько ярких линий
    \begin{scope}
      \clip (current page.south west) rectangle (current page.north east);
      % частая мелкая сетка
      \foreach \i in {0,...,40}{
        \draw[coverline, opacity=0.10, line width=0.3pt]
          ($(current page.north west)+(-3,-\i*0.8)$) --
          ($(current page.north east)+(0,-\i*0.8-13)$);}
      % дуги из верхнего правого угла
      \foreach \r in {5,7,...,29}{
        \draw[coverline, opacity=0.16, line width=0.5pt]
          ($(current page.north east)+(-1,1)$) circle (\r);}
      % встречные диагонали снизу
      \foreach \i in {0,...,14}{
        \draw[coverline, opacity=0.09, line width=0.4pt]
          ($(current page.south west)+(0,\i*1.6)$) --
          ($(current page.south east)+(0,\i*1.6+11)$);}
      % акцентные линии
      \foreach \i in {4,11,18,25}{
        \draw[coverline, opacity=0.34, line width=0.6pt]
          ($(current page.north west)+(-3,-\i*0.8)$) --
          ($(current page.north east)+(0,-\i*0.8-13)$);}
      \foreach \r in {13,21}{
        \draw[coverline, opacity=0.30, line width=0.7pt]
          ($(current page.north east)+(-1,1)$) circle (\r);}
    \end{scope}

    \draw[coverline, opacity=0.45, line width=0.7pt]
      ($(current page.north west)+(2cm,-2.6cm)$) --
      ($(current page.north east)+(-2cm,-2.6cm)$);

    % верхняя строка
    \node[anchor=north west, white, inner sep=0pt]
      at ($(current page.north west)+(2cm,-1.75cm)$)
      {\fontsize{11}{13}\selectfont\bfseries\MakeUppercase{\coursemeta}};

    % заголовок
    \node[anchor=west, white, inner sep=0pt, text width=15.5cm, align=left]
      at ($(current page.west)+(2cm,1.1cm)$)
      {{\fontsize{44}{50}\bfseries\selectfont \subjectname\par}
       \vspace{0.55cm}
       {\fontsize{17}{20}\selectfont\color{white!72} Конспект лекций\par}};

    % лектор
    \draw[coverline, opacity=0.40, line width=0.6pt]
      ($(current page.west)+(2cm,-3.55cm)$) -- ($(current page.west)+(8cm,-3.55cm)$);
    \node[anchor=north west, white, inner sep=0pt, text width=12cm, align=left]
      at ($(current page.west)+(2cm,-4.15cm)$)
      {{\fontsize{10}{12}\selectfont\color{white!55}\MakeUppercase{Лектор}\par}
       \vspace{2pt}
       {\fontsize{15}{18}\selectfont \lecturer\par}};

    % низ
    \node[anchor=south west, white, inner sep=0pt]
      at ($(current page.south west)+(2cm,1.9cm)$)
      {\fontsize{11}{13}\selectfont\color{white!55} Санкт-Петербург,~\the\year};
  \end{tikzpicture}
  \endgroup
  \end{titlepage}
  \clearpage
}

% ---------- страница содержания -----------------------------------
\newcommand{\maketoc}{%
  \thispagestyle{empty}
  \begingroup
  \setcounter{tocdepth}{1}
  \tableofcontents
  \vfill
  \begingroup\footnotesize\sffamily\color{inkSoft}
  \noindent{\color{ruleC}\rule{\textwidth}{0.5pt}}\\[6pt]
  \noindent Составитель: \compiler\ \ \href{https://github.com/\compilergh}{@\compilergh}\\[3pt]
  \noindent\addedglyph\ \ фрагмент дописан составителем: на лекции не приводился
  или был дан устно.\par
  \endgroup
  \endgroup
  \clearpage
}
